\scanpage{105}
\mspage{3}
% archive foliation: 104
Sections de $P-\sigma(X)$ = connexions sur $\omega$
\[
=\text{op. [diff. d'ordre 1] }\underline\omega\longrightarrow\underline\omega^{\unsure{2}}\deleted{\unsure{\omega}}
\]
satisfaisant $D(f\omega)=fD\omega+\omega\,df$
\unsure{forme} torseur sous $\underline\omega$

À toute telle section correspond une connexion projective particulière $c_D$ --- l'unique conn. projective pour laquelle cette section est horizontale.
Si on se fixe une section du \inserted{$P-\sigma(X)$, i.e. une section du} torseur des connexions sur $\underline\omega$ \inserted{\unsure{i.e. si on a une connexion} $D_1:\omega\longrightarrow\omega^2$}, i.e. un splitting de $\mathcal L$ (ou encore de $\mathcal H$\ldots), soit $D_0$\deleted{\unsure{D}}, \deleted{ce qui donne lieu à \unsure{une appl. de}}
\[
\boxed{c_{D_0+\underbrace{\underline\omega_1}_{\text{entrée}}}=\deleted{\unsure{D_0(\omega_1)+\omega_1^2}}c_{D_0}+2D_0(\omega_1)+\omega_1^2}
\]
\unsure{Et vérifions}
\[
\begin{aligned}
c_{D_0+\omega_1+\bar\omega_1}
&=\underbrace{c_{D_0+\omega_1}}_{c_{D_0}+\unsure{D_0\omega_1}+\omega_1^2}
+\underbrace{(D_0+\omega_1)(\bar\omega_1)+\bar\omega_1^{\unsure{2|1}}}_{2D_0(\bar\omega_1)+2\omega_1\bar\omega_1+\bar\omega_1^2}\\
&=c_{D_0}+2D_0(\omega_1+\bar\omega_1)+\omega_1^2+\bar\omega_1^2+2\omega_1\bar\omega_1\\
&=c_{D_0}+2D_0(\omega_1\unsure{+|,}\bar\omega_1)+(\omega_1+\bar\omega_1)^2\qquad\text{OK}
\end{aligned}
\]
Équation d'horizontalité de $D_0+\omega_1$ pour une connexion proj. $c_{D_0}+\bar\omega_2$ :
\[
c_{D_0}+\bar\omega_2=c_{D_0+\omega_1}\quad\text{i.e.}
\]
\[
\boxed{2D_0(\omega_1)+\omega_1^2=\bar\omega_2}
\]
\unsure{Considérons} la connexion
\[
D_t(\lambda\,dt)\longrightarrow\frac{d\lambda}{dt}dt^2\quad\text{i.e.}
\]
Ou :
\[
D_t(\omega)=\left(\frac d{dt}\left(\frac\omega{dt}\right)\right)dt^2
\]
\[
\deleted{D_0(f\lambda\,dt)}=\left[d\left(\frac\omega{dt}\right)\right]dt
\]
Ou bien
\deleted{$D_t(f\omega)=fD_t(\omega)+{}$}

\margin{\diagram{\(\begin{array}{r|l|l|l}
&\underline\omega&\underline\omega^0&\underline\omega^{-1}\\
\deleted{\unsure{\omega+}}&&&\\
\alpha_1+\alpha_{-1}=a&&&\\
\rho_a\omega_1=& &\deleted{\omega_1\alpha_{-1}}\inserted{\alpha_{-1}\omega_1}&\deleted{\unsure{\alpha_{-1}\omega_{-1}}}\\
\rho_a\lambda=&\lambda\alpha_1&&+\lambda\alpha_{-1}\\
\rho_a\omega_{-1}=&&\alpha_1\omega_{-1}&\\
(\rho_a)^2\omega_1=&(\alpha_{-1}\alpha_1)\omega_1&&+(\alpha_{-1})^2\omega_{\unsure{1|-1}}\\
(\rho_a)^2\lambda=&&2\alpha_1\alpha_{-1}\lambda&\\
(\rho_a)^2\omega_{-1}=&(\alpha_1^2)\omega_{-1}&&+(\alpha_1\alpha_{-1})\omega_{-1}\\
(\rho_a)^3\omega_1=&&2\alpha_{-1}^2\alpha_1\omega_1&\\
(\rho_a)^3\lambda=&2\alpha_1^2\alpha_{-1}\lambda&&+2\alpha_1\alpha_{-1}^2\lambda
\end{array}\)}}
\margin{TSVP}
